\chapter{Problem Definition and Proposed Algorithm}
This chapter formalizes the Routing and Spectrum Allocation (RSA) challenge within multi-graph optical topologies and presents the core algorithmic solutions developed for this thesis. It details the methodologies used to model physical networks with parallel links, compute required bandwidth and slots, discover optimal and alternative paths, and dynamically allocate spectrum using a novel bitmap intersection approach.
\section{Formal Problem Statement}
To address the complexities of lightpath provisioning in modern elastic optical networks, the physical topology must first be formalized into a mathematical model. The optical network is represented as a directed multi-graph, denoted as $G=(N,L)$. In this model, $N$ represents the set of optical network nodes (e.g., transponders, muxponders, and ROADMs), and $L$ represents the set of bidirectional physical optical links interconnecting these nodes. A defining characteristic of the targeted architecture is the presence of parallel links—multiple distinct physical fibers connecting the same node pair $(u,v)\in N$. Therefore, the edge set $L$ is defined such that multiple edges can exist between any two vertices, formally expressed as a multi-graph where a link $l\in L$ is defined by the tuple $(u,v,k)$, with $k$ serving as a unique identifier for a specific parallel fiber between nodes $u$ and $v$. Each link $l$ supports a total spectrum capacity divided into a set of frequency slots, $W$, adhering to standard ITU-T flexible grid specifications.

The network is governed by a centralized control plane, orchestrated by the TeraFlowSDN ecosystem. Within this environment, it is assumed that several active lightpaths, denoted as $LP_i(N_s, N_d)$, already exist, utilizing specific continuous and contiguous blocks of spectrum $w_i\subset W$ across various links to connect source nodes $N_s$ to destination nodes $N_d$. The core challenge arises when a new lightpath request is submitted to the controller. The controller cannot simply compute the shortest path by treating the network as a standard simple graph. It must account for the current spectral state of the network, recognizing that certain links may be nearing capacity. Most critically, the routing algorithm must resolve the stringent physical constraints imposed by parallel links terminating at the same ROADM device. Specifically, in modern deployments, multiple parallel fibers originating from the same muxponder are frequently directed to the same Shared Risk Group (SRG) and subsequently multiplexed into the same Degree (DEG) of a ROADM. Because these signals share the same physical amplification and switching infrastructure within the DEG, they are subject to strict spectral isolation constraints. If active lightpaths currently utilize a specific set of frequency slots on link $l_(u,v,1)$, those identical frequencies cannot be reused on a parallel link $l_(u,v,2)$ connecting the same devices, as doing so would result in severe optical interference. Therefore, when establishing a new lightpath across a parallel edge $l_(u,v,k)$, the routing logic cannot evaluate the spectrum availability of link $k$ in isolation. It must perform a concurrent evaluation of the entire spatial bundle of parallel links connecting $u$ and $v$ to ensure that the newly allocated contiguous spectrum block does not overlap with any frequencies currently active on sibling links $i$ or $j$. The formal objective of the \texttt{ParallelOpticalController} is to efficiently compute the optimal routing path and allocate the necessary contiguous spectrum for incoming requests, specifically navigating the multi-dimensional constraints of parallel links, while minimizing overall network blocking probability and computational latency.
\section{Graph Modelling with Parallel Links (topologies)}
To efficiently navigate the multi-dimensional constraints established in the formal problem statement, the \texttt{ParallelOpticalController} avoids performing complex computations directly on the raw physical topology. Instead, the topological graph is logically segregated into specialized, abstracted versions. Acting as the centralized intelligence of the network, the TeraFlowSDN ecosystem provides the real-time database and telemetry required to construct, filter, and dynamically update these graphical models prior to processing any lightpath request. The core topology is first logically abstracted into an available-resource graph, denoted as $G_{free}$. The purpose of $G_{free}$ is to immediately prune physical resources that cannot mathematically support a new connection, thereby shrinking the search space. During its construction, the controller queries the TeraFlowSDN database and actively excludes fully occupied Optical Multiplex Section (OMS) links between ROADMs, as well as currently utilized Optical Channel (OCH) ports between muxponders and ROADMs. The resulting $G_{free}$ remains a multi-graph containing parallel links, but it strictly represents the network's actual residual capacity. While $G_{free}$ accurately reflects available resources, executing shortest-path algorithms directly on a dense multi-graph introduces significant, often exponential, computational overhead. To bypass this state-space explosion, $G_{free}$ is further abstracted into a simplified graph, denoted as $G_{simple}$. In this reduction, all parallel physical fibers connecting any adjacent node pair (u,v) are collapsed into a single representative logical edge. This transformation allows the controller to rapidly execute standard shortest-path and alternative-path computations (such as Dijkstra's shortest path algorithm) to discover the optimal topological routes from the source to the destination. Crucially, $G_{simple}$ is used strictly for topological path discovery, not for spectrum allocation. Once the candidate node-to-node routes are identified, the controller projects these paths back onto the multi-graph $G_{free}$ to perform the actual Routing and Spectrum Allocation (RSA). By evaluating the physical routes against $G_{free}$, the controller natively re-introduces the spatial domain—allowing it to concurrently evaluate all parallel fibers for a given hop. This dual-graph methodology ensures that the path finding phase remains computationally lightweight, while the final RSA phase mathematically guarantees that utilized spectrum is accurately nullified and the contiguity constraints are reliably satisfied across parallel links.
\section{Proposed RSA Algorithm}
The proposed Routing and Spectrum Allocation (RSA) algorithm operates through a sequence of integrated stages to dynamically provision lightpaths. The process is initiated when a service request is submitted via the TeraFlowSDN WebUI using a JSON descriptor. The controller parses this descriptor to extract critical constraints, such as the requested bandwidth and the hop-count limit for alternative path computation. Following this extraction, the algorithm systematically executes bandwidth and slot calculation, topological path discovery, and spatial spectrum allocation analysis across parallel links. Finally, rather than forcing an automated assignment, the architecture presents the computed valid path permutations intuitively via the control plane, allowing the network operator to explicitly select the desired spatial route for final spectrum acquisition and assignment.
\begin{figure}[htbp]
    \centering
    \includegraphics[width=1\textwidth]{full_algo_er.png}
    \caption{RSA operation using ParallelOpticalController in TeraFlowSDN}
    \label{fig:full_algo_er}
\end{figure}
\subsection{Bandwidth \& Slot Computation}
The initial phase of the proposed algorithm determines the required spectral bandwidth ($B_r$) for the requested optical service. Upon receiving a connection request, the controller extracts the target transmission bitrate ($R_b$), expressed in Gigabits per second (Gbps), directly from the submitted TeraFlowSDN service descriptor. The required bandwidth is then mathematically derived using the following formulation:
\[
    B_r = \frac{R_b}{n}(1 + \alpha)
\]
where $n=log_2(M)$, and M represents the specific modulation format utilized. While several factors influence the required bandwidth calculation \cite{ntiaBandwidth}, this standard formula simplifies the process for the controller.

Once the required bandwidth $B_r$ is established, the controller computes the discrete number of frequency slots ($N_r$) needed to accommodate the lightpath. In accordance with ITU-T flexible grid standards, while 12.5 GHz typically represents the minimum routable spectrum width, the nominal central frequency granularity of 6.25 GHz is utilized as the fundamental unit for this computation. This finer 6.25 GHz granularity is deliberately adopted to streamline the mathematical modeling and facilitate the subsequent bitmap-based spectrum assignment logic. The required number of slots is calculated as follows:
\[
    N_r = \frac{B_r}{6.25}
\]
Upon completion of this phase, the computed parameters, $B_r$ and $N_r$, are temporarily cached within the controller's operational memory. These values serve as the primary spectral constraints for the ensuing topological path discovery and spatial spectrum allocation phases.
\subsection{Path Discovery}
Following the determination of the required bandwidth and frequency slots, the controller advances to the topological path discovery phase. To perform fast and computationally efficient routing, the algorithm leverages NetworkX\cite{networkX}, a robust and highly optimized Python library designed for complex network analysis. During this phase, the routing engine queries the simplified graph abstraction, $G_{simple}$—where parallel physical fibers are collapsed into single logical edges. By routing over $G_{simple}$ rather than the raw physical multi-graph, the controller effectively prevents the state-space explosion and latency bottlenecks that plague traditional routing algorithms. The path computation engine is designed to execute two distinct routing operations to provide the user with maximum flexibility: it first identifies the absolute shortest path, and subsequently computes a set of viable alternative routes. This alternative path computation is governed by a user-defined hop-relaxation constraint, formulated as $H_{shortest}+n$, where $H_{shortest}$ represents the hop count of the primary shortest path and $n$ is a positive integer configured via the TeraFlowSDN service descriptor. The algorithm filters the topology to discover and validate alternative routes that strictly conform to this exact hop-count limit.

Once these node-level routes are successfully discovered in $G_{simple}$, the controller projects them back onto the available-resource multi-graph, $G_{free}$. Naturally, in a multi-fiber environment, a single node-to-node logical route expands into multiple physical path permutations between the source $(N_s)$ and destination $(N_d)$ devices. For standard operational deployment within the TeraFlowSDN ecosystem, the controller is configured to default to the first valid physical path object among these parallel permutations that satisfies the spectral constraints. However, to comprehensively evaluate the controller's spatial awareness, several alternate spatial selection strategies have been investigated, with their statistical impacts on network performance and blocking probability detailed in the subsequent analytical sections of this thesis.
\subsection{Bitmap-Based Spectrum Allocation}
Upon successful path discovery, the controller initiates the spectrum allocation phase using the primary shortest-path object derived from Dijkstra's algorithm, along with the previously computed slot requirement, $N_r$. While the controller defaults to evaluating the primary shortest path, RSA computation for any discovered alternative routes can be executed subsequently at the operator's discretion. The allocation process begins by querying the controller's internal cache to retrieve the real-time operational state of all physical links comprising the selected path object. The foundational step in this RSA methodology is establishing a uniform spectral reference framework. The controller analyzes the path to identify the maximum spectral band supported by any traversed device. Within the cache, the current spectrum utilization of each port is stored not just as discrete slots, but mathematically encoded as an integer value representing a binary sequence. This integer, referred to as the \texttt{bitmap-value}, is calculated using the following formulation:
\[
    \text{bitmap-value} = 2^N - 1
\]
The integer representing the widest supported operational band (typically the C+L band) is designated as the \texttt{bitmap\_initial}. This value establishes the maximum bit-length reference for all subsequent logical operations. A defining feature of this RSA algorithm is its native, unconditional sensitivity to multi-fiber topologies. Regardless of whether a specific hop explicitly utilizes multiple fibers, the algorithm inherently executes a parallelism check, ensuring that spectrum continuity and contiguity are dynamically evaluated across all available spatial dimensions between adjacent nodes. For a given hop along the path from source $N_s$ to destination $N_d$, the controller aggregates all available source ports into a set $P_{src}$ and all corresponding destination ports into a set $P_{dst}$. The current spectrum utilization of each port is retrieved as its integer \texttt{Bitmap\_value} and converted into its binary representation. To ensure mathematical alignment, each binary sequence is zero-padded or expanded to perfectly match the bit-length of the reference \texttt{bitmap\_initial}. Once aligned, the controller executes a concurrent bitwise intersection (logical AND operations) across the expanded bitmaps of all eligible source and destination ports within that hop. This concurrent spatial evaluation collapses the multi-fiber complexity into a single, unified mathematical representation, yielding a definitive \texttt{bitmap\_final}. This \texttt{bitmap\_final} sequence explicitly maps the continuous and contiguous spectrum gaps that are simultaneously available across the entire parallel link bundle. The controller then applies a standard First-Fit policy to the \texttt{bitmap\_final} sequence to identify the lowest-indexed contiguous block of zeros (available spectrum) that meets the required length $N_r$. The results of this deterministic spatial-spectral evaluation, including the final computed bitmap and the corresponding physical path hops, are subsequently relayed to the TeraFlowSDN WebUI for operator review.
\subsection{Slot Acquisition \& Assignment}
The final phase of the proposed RSA algorithm transitions from computation to operational execution, facilitating the actual acquisition and assignment of spectral slots. Rather than autonomously enforcing an allocation, the architecture employs an interactive, operator-in-the-loop paradigm. The network operator is presented with the computed primary and alternative path objects—along with their corresponding available spectrum blocks—and is granted the flexibility to explore these options and explicitly select the preferred spatial route via the TeraFlowSDN WebUI. Upon operator's confirmation, the controller initiates the slot acquisition protocol for the selected path. It retrieves the current link information for the chosen route and updates the spectral state by converting the newly allocated continuous blocks from their binary bitmap representations back into integer values. These updated integer values are subsequently pushed to the centralized database via the TeraFlowSDN Context component (e.g., using the Context client API) to accurately reflect the network's newly diminished capacity.

Concurrently, the controller executes a series of state updates across the topology. At the central database level, the specific spatial device ports utilized by the lightpath are flagged as \texttt{used=True} to prevent future routing collisions. At the orchestration level, the computed lightpath object is transitioned to an "active" state, and the overarching service request is formally marked as "ACTIVE" within the central knowledge base. Operators can subsequently monitor these allocated resources, port usages, and path details through the WebUI. Crucially, within the current scope of this implementation, the \texttt{ParallelOpticalController} operates strictly as a control-plane and orchestration entity. While it comprehensively updates the centralized database and topological state to enforce resource locking, it does not currently generate or transmit southbound physical configuration commands (such as NETCONF or RESTCONF) to alter the hardware interfaces of the physical optical devices. All spectral acquisition and state management are logically enforced at the controller level.
\section{Complexity Analysis}
The computational efficiency of the proposed \texttt{ParallelOpticalController} is fundamentally anchored in its ability to decouple topological routing from spatial spectrum allocation. The initialization phase, which constructs the abstracted multi-graph $(G_{free})$ and simplified routing graph $(G_{simple})$, traverses the network state in strictly linear time, $O(|N|+|L|)$—effectively $O(n)$ relative to the total topology size ($n$ represents a unit consisted of $N,L$)—ensuring predictable memory and processing scaling as physical devices and links are added. Subsequently, by executing topological path discovery exclusively on the collapsed graph $G_{simple}$, the controller avoids the state-space explosion typical of multi-graph environments, maintaining Dijkstra's standard log-linear time complexity of $O(|L_{simple}|+|L \times logL|)$ (where $L_{simple}$ represents the links involved in the shortest-path). The most significant computational breakthrough, however, emerges during the Routing and Spectrum Allocation (RSA) phase.

While traditional SDN implementations evaluate multi-fiber topologies sequentially—resulting in an exponential time complexity of $O(L_{parallel}^H)$, where $L_{parallel}$ represents the number of parallel links per hop and $H$ denotes the path length(measured in hops)—the proposed algorithm evaluates spatial bundles concurrently. Because the bitwise intersection mathematical logic resolves spectral contiguity across all parallel ports in constant O(1) time at the CPU hardware level, the overall time complexity of the RSA phase is strictly bounded by $O(L_{parallel}\times H)$. This additive, linear scaling with respect to path length and fiber density, coupled with a minimal spatial footprint of $O(N+L)$ for graph storage and negligible auxiliary memory for the binary string operations, mathematically guarantees that the controller can provision complex, multi-fiber lightpaths without succumbing to the exponential latency bottlenecks inherent in conventional routing algorithms. To empirically validate these theoretical claims, comprehensive performance tests were conducted across varying network topologies, including well-known reference networks such as the Pan-European (PANEU) and NSFNet models. As demonstrated in the results below, the controller successfully maintains the mathematically established time and space complexities, exhibiting highly efficient scaling even when multi-fiber parallel links are introduced. Ultimately, the convergence of theoretical analysis and empirical performance metrics confirms the robust computational efficiency of the \texttt{ParallelOpticalController}. By leveraging concurrent spatial evaluation to eliminate sequential processing bottlenecks, the architecture ensures highly scalable, real-time lightpath provisioning capable of meeting the rigorous demands of next-generation, multi-fiber elastic optical networks.